% !TEX root = ../main.tex
%% Devoir maison n° 4 — Inégalités de concentration (chapitres 6, 11, 17)
%% Résultats vérifiés par verifications/eval_dm.py.

\sujetEval{DM4}{MPSI · Devoir maison}{Devoir maison \numero4}
  {À rendre en une semaine\\Inégalités de concentration\\Chapitres 6, 11 et 17}{3}
  {Partie A/L'inégalité de Bienaymé-Tchebychev/Q1 et Q2,Partie B/L'inégalité de Hoeffding/Q3 à Q6,%
   Partie C/Application à un sondage/Q7 à Q9}
\enParties
\consigne{On rappelle que $\operatorname{ch}t=\frac{\mathrm{e}^{t}+\mathrm{e}^{-t}}{2}$
et $\operatorname{th}t=\frac{\operatorname{sh}t}{\operatorname{ch}t}$. Les
valeurs numériques demandées en partie C se calculent à la main, avec
$\ln40\approx3{,}689$.}

Soit $n\geqslant1$ et $X_1,\dots,X_n$ des variables aléatoires
indépendantes, définies sur un même espace probabilisé fini, de même loi
\[P(X_i=1)=P(X_i=-1)=\frac12.\]
On pose $S_n=X_1+\dots+X_n$ : c'est la position, après $n$ pas, d'un
marcheur qui avance ou recule d'un pas au hasard. On cherche à majorer la
probabilité que $S_n$ s'écarte beaucoup de $0$.

\begin{exoEval}[L'inégalité de Bienaymé-Tchebychev]{}{}
\begin{questions}
  \item Calculer $E(S_n)$ et $V(S_n)$.
  \item En déduire que pour tout $a>0$ :
    $P\bigl(|S_n|\geqslant a\bigr)\leqslant\dfrac{n}{a^{2}}$.
\end{questions}
\end{exoEval}

\begin{exoEval}[L'inégalité de Hoeffding]{}{}
Soit $a>0$.
\begin{questions}[start=3]
  \item Soit $t>0$. À l'aide de l'inégalité de Markov appliquée à
    $\mathrm{e}^{tS_n}$, montrer que
    $P(S_n\geqslant a)\leqslant\mathrm{e}^{-ta}\,E\bigl(\mathrm{e}^{tS_n}\bigr)$.
  \item Montrer que $E\bigl(\mathrm{e}^{tS_n}\bigr)=(\operatorname{ch}t)^{n}$.
  \item Étudier sur $\R_{+}$ la fonction
    $g:t\mapsto\frac{t^{2}}{2}-\ln(\operatorname{ch}t)$, et en déduire que
    $\operatorname{ch}t\leqslant\mathrm{e}^{t^{2}/2}$ pour tout réel $t$.
  \item En déduire que $P(S_n\geqslant a)\leqslant\exp\bigl(-ta+\frac{nt^{2}}{2}\bigr)$
    pour tout $t>0$, puis, en choisissant $t$ au mieux,
    \[P(S_n\geqslant a)\leqslant\exp\Bigl(-\frac{a^{2}}{2n}\Bigr)\qquad\text{et}\qquad
      P\bigl(|S_n|\geqslant a\bigr)\leqslant2\exp\Bigl(-\frac{a^{2}}{2n}\Bigr).\]
\end{questions}
\end{exoEval}

\begin{exoEval}[Application à un sondage]{}{}
On lance $n$ fois une pièce équilibrée, $X_i$ valant $1$ si le $i$-ième
lancer donne « pile » et $-1$ sinon. On note $F_n$ la fréquence de « pile »
et l'on fixe $\varepsilon>0$.
\begin{questions}[start=7]
  \item Montrer que $F_n=\dfrac{n+S_n}{2n}$, puis que
    \[P\Bigl(\Bigl|F_n-\frac12\Bigr|\geqslant\varepsilon\Bigr)\leqslant\frac{1}{4n\varepsilon^{2}}
      \qquad\text{et}\qquad
      P\Bigl(\Bigl|F_n-\frac12\Bigr|\geqslant\varepsilon\Bigr)\leqslant2\,\mathrm{e}^{-2n\varepsilon^{2}}.\]
  \item On prend $\varepsilon=0{,}05$. Pour chacune des deux majorations,
    déterminer le plus petit $n$ qui garantit
    $P\bigl(|F_n-\frac12|\geqslant0{,}05\bigr)\leqslant0{,}05$.
  \item Comparer les deux majorations lorsque $n\to+\infty$. Commenter.
\end{questions}
\end{exoEval}
